\documentclass[11pt]{article}
\usepackage{mathunal-base}
\mathunalfuente{serif}
\providecommand{\nota}[1]{\par\smallskip\noindent{\small\itshape\color{unalblue}#1}\par\smallskip}
\newlist{opciones}{enumerate}{1}
\setlist[opciones]{label=(\roman*),itemsep=1pt,topsep=2pt,leftmargin=1.8em,font=\normalfont}
\newcommand{\rp}{\underline{\hspace{3cm}}}

\renewcommand{\examtitulo}{Segundo Parcial de Ecuaciones Diferenciales}
\renewcommand{\examfecha}{3 de Junio de 2023}

\begin{document}

\examheader

\vspace{4pt}
\noindent\textit{Duración: 1 hora y 50 minutos. El examen consta de cuatro preguntas en dos hojas impresas por ambos lados; antes de empezar verifique que su examen esté completo y que sus datos sean correctos. No desengrape su examen ni añada otras grapas o su examen será anulado. En las preguntas con procedimiento justifique sus respuestas en los espacios asignados. En los puntos 1 y 2 debe justificar paso a paso su procedimiento; en los puntos 3 y 4 sólo se calificará la respuesta, no es necesario que escriba el procedimiento. No está permitido hacer preguntas, usar calculadoras, sacar hojas en blanco ni ningún tipo de apuntes durante el examen. Por favor verifique que su celular esté apagado.}

\vspace{6pt}
\noindent\textbf{1. (35\%)} Usando el método de variación de parámetros halle la solución general en $(-\infty,\infty)$ de la ecuación diferencial
\[
y''-5y'+6y=2e^{2x}.
\]

\vspace{6cm}

\newpage

\noindent\textbf{2. (10\%)} Suponga que $p$ y $q$ son funciones continuas y que $y_1$ y $y_2$ son soluciones de la ecuación diferencial
\[
y''+p(x)y'+q(x)y=0
\]
sobre un intervalo abierto $I$. Demuestre que si $y_1$ y $y_2$ alcanzan su valor máximo o mínimo en el mismo punto $x_0\in I$, entonces sobre ese intervalo no pueden ser un conjunto fundamental de soluciones de tal ecuación.

\vspace{6cm}

\newpage

\noindent\textbf{3. (25\%)} En este ejercicio sólo se calificará respuesta, no es necesario que escriba el procedimiento.

\vspace{0.2cm}
\noindent\textbf{a)} Una masa que pesa 32 libras alarga 8 pies un resorte. Al inicio la masa se libera desde un punto que está $\frac{1}{2}$ pie debajo de la posición de equilibrio con una velocidad de 1 pie/seg hacia arriba. En los siguientes numerales complete según corresponda.
\begin{opciones}
\item \textbf{(10\%)} Al escribir la posición en la forma $x(t)=A\sin(2t+\Phi)$, los valores de $A$ y $\Phi$ son $A=$~\rp\ y $\Phi=$~\rp
\item \textbf{(5\%)} Los ciclos completos que habrá realizado la masa al final de $5\pi$ segundos son \rp\ ciclos.
\end{opciones}

\vspace{0.3cm}
\noindent\textbf{b) (10\%)} La posición de un cuerpo, cuya masa es 1 slug, que está sujeto a un resorte, está dada por
\[
x(t)=e^{-2t}\sin\!\left(2t+\frac{3\pi}{4}\right)
\]
pies, para cada instante $t\geq 0$ (medido en segundos).
\begin{opciones}
\item \textbf{(5\%)} La posición inicial de la masa es $x(0)=$~\rp\ pies y su velocidad inicial es $x'(0)=$~\rp\ pies/seg. (Sus respuestas deben ser números explícitos).
\item \textbf{(5\%)} En el instante $t=3$ seg, determine la posición de la masa e indique hacia dónde se dirige.
\end{opciones}

\newpage

\noindent\textbf{4. (30\%)} En los literales a), b), c) y d) complete según corresponda. Sólo se calificará respuesta, no es necesario que escriba el procedimiento.

\vspace{0.25cm}
\noindent\textbf{a) (8\%)} Si $x^{-2}$ y $x^{-2}\ln x$ son soluciones de una ecuación diferencial de Cauchy-Euler homogénea de orden 2 en $(0,\infty)$, entonces dicha ecuación es: \rp

\vspace{0.3cm}
\noindent\textbf{b) (10\%)} Considere la ecuación diferencial
\[
y^{(4)}+8y'''+16y''=-xe^{x}+7.
\]
\begin{opciones}
\item \textbf{(4\%)} Un conjunto fundamental de soluciones de la ecuación diferencial homogénea asociada es \rp
\item \textbf{(6\%)} Una forma apropiada de solución particular, según el método de coeficientes indeterminados, para la ecuación diferencial es: \rp\ (No calcule las constantes, sólo escriba la forma adecuada).
\end{opciones}

\vspace{0.3cm}
\noindent\textbf{c) (8\%)} Considere la ecuación diferencial homogénea con coeficientes constantes
\[
y''-y'+qy=0, \qquad x\in\mathbb{R}.
\]
Los valores de $q$ para que la solución general de esta ecuación diferencial sea de la forma $y=c_1e^{r_1x}+c_2e^{r_2x}$, con $r_1$ y $r_2$ en $\mathbb{R}$ y $r_1\neq r_2$, son: \rp

\vspace{0.3cm}
\noindent\textbf{d) (4\%)} La posición $x(t)$ de cierto sistema masa-resorte satisface el problema de valor inicial
\[
\frac{9}{25}x''+16x=2\sin(\gamma t), \qquad x(0)=0,\ x'(0)=0.
\]
El valor de $\gamma$ que hace que ocurra el fenómeno de resonancia es: \rp

\vfill
\rule{\linewidth}{0.4pt}\\[1pt]
{\scriptsize\textit{Transcripción digital --- MathUNAL}\hfill\textcolor{unalblue}{\textbf{1}}}

\end{document}
